% Options for packages loaded elsewhere
\PassOptionsToPackage{unicode}{hyperref}
\PassOptionsToPackage{hyphens}{url}
%
\documentclass[
]{article}
\usepackage{amsmath,amssymb}
\usepackage{iftex}
\ifPDFTeX
  \usepackage[T1]{fontenc}
  \usepackage[utf8]{inputenc}
  \usepackage{textcomp} % provide euro and other symbols
\else % if luatex or xetex
  \usepackage{unicode-math} % this also loads fontspec
  \defaultfontfeatures{Scale=MatchLowercase}
  \defaultfontfeatures[\rmfamily]{Ligatures=TeX,Scale=1}
\fi
\usepackage{lmodern}
\ifPDFTeX\else
  % xetex/luatex font selection
\fi
% Use upquote if available, for straight quotes in verbatim environments
\IfFileExists{upquote.sty}{\usepackage{upquote}}{}
\IfFileExists{microtype.sty}{% use microtype if available
  \usepackage[]{microtype}
  \UseMicrotypeSet[protrusion]{basicmath} % disable protrusion for tt fonts
}{}
\makeatletter
\@ifundefined{KOMAClassName}{% if non-KOMA class
  \IfFileExists{parskip.sty}{%
    \usepackage{parskip}
  }{% else
    \setlength{\parindent}{0pt}
    \setlength{\parskip}{6pt plus 2pt minus 1pt}}
}{% if KOMA class
  \KOMAoptions{parskip=half}}
\makeatother
\usepackage{xcolor}
\usepackage[margin=1in]{geometry}
\usepackage{color}
\usepackage{fancyvrb}
\newcommand{\VerbBar}{|}
\newcommand{\VERB}{\Verb[commandchars=\\\{\}]}
\DefineVerbatimEnvironment{Highlighting}{Verbatim}{commandchars=\\\{\}}
% Add ',fontsize=\small' for more characters per line
\usepackage{framed}
\definecolor{shadecolor}{RGB}{248,248,248}
\newenvironment{Shaded}{\begin{snugshade}}{\end{snugshade}}
\newcommand{\AlertTok}[1]{\textcolor[rgb]{0.94,0.16,0.16}{#1}}
\newcommand{\AnnotationTok}[1]{\textcolor[rgb]{0.56,0.35,0.01}{\textbf{\textit{#1}}}}
\newcommand{\AttributeTok}[1]{\textcolor[rgb]{0.13,0.29,0.53}{#1}}
\newcommand{\BaseNTok}[1]{\textcolor[rgb]{0.00,0.00,0.81}{#1}}
\newcommand{\BuiltInTok}[1]{#1}
\newcommand{\CharTok}[1]{\textcolor[rgb]{0.31,0.60,0.02}{#1}}
\newcommand{\CommentTok}[1]{\textcolor[rgb]{0.56,0.35,0.01}{\textit{#1}}}
\newcommand{\CommentVarTok}[1]{\textcolor[rgb]{0.56,0.35,0.01}{\textbf{\textit{#1}}}}
\newcommand{\ConstantTok}[1]{\textcolor[rgb]{0.56,0.35,0.01}{#1}}
\newcommand{\ControlFlowTok}[1]{\textcolor[rgb]{0.13,0.29,0.53}{\textbf{#1}}}
\newcommand{\DataTypeTok}[1]{\textcolor[rgb]{0.13,0.29,0.53}{#1}}
\newcommand{\DecValTok}[1]{\textcolor[rgb]{0.00,0.00,0.81}{#1}}
\newcommand{\DocumentationTok}[1]{\textcolor[rgb]{0.56,0.35,0.01}{\textbf{\textit{#1}}}}
\newcommand{\ErrorTok}[1]{\textcolor[rgb]{0.64,0.00,0.00}{\textbf{#1}}}
\newcommand{\ExtensionTok}[1]{#1}
\newcommand{\FloatTok}[1]{\textcolor[rgb]{0.00,0.00,0.81}{#1}}
\newcommand{\FunctionTok}[1]{\textcolor[rgb]{0.13,0.29,0.53}{\textbf{#1}}}
\newcommand{\ImportTok}[1]{#1}
\newcommand{\InformationTok}[1]{\textcolor[rgb]{0.56,0.35,0.01}{\textbf{\textit{#1}}}}
\newcommand{\KeywordTok}[1]{\textcolor[rgb]{0.13,0.29,0.53}{\textbf{#1}}}
\newcommand{\NormalTok}[1]{#1}
\newcommand{\OperatorTok}[1]{\textcolor[rgb]{0.81,0.36,0.00}{\textbf{#1}}}
\newcommand{\OtherTok}[1]{\textcolor[rgb]{0.56,0.35,0.01}{#1}}
\newcommand{\PreprocessorTok}[1]{\textcolor[rgb]{0.56,0.35,0.01}{\textit{#1}}}
\newcommand{\RegionMarkerTok}[1]{#1}
\newcommand{\SpecialCharTok}[1]{\textcolor[rgb]{0.81,0.36,0.00}{\textbf{#1}}}
\newcommand{\SpecialStringTok}[1]{\textcolor[rgb]{0.31,0.60,0.02}{#1}}
\newcommand{\StringTok}[1]{\textcolor[rgb]{0.31,0.60,0.02}{#1}}
\newcommand{\VariableTok}[1]{\textcolor[rgb]{0.00,0.00,0.00}{#1}}
\newcommand{\VerbatimStringTok}[1]{\textcolor[rgb]{0.31,0.60,0.02}{#1}}
\newcommand{\WarningTok}[1]{\textcolor[rgb]{0.56,0.35,0.01}{\textbf{\textit{#1}}}}
\usepackage{graphicx}
\makeatletter
\def\maxwidth{\ifdim\Gin@nat@width>\linewidth\linewidth\else\Gin@nat@width\fi}
\def\maxheight{\ifdim\Gin@nat@height>\textheight\textheight\else\Gin@nat@height\fi}
\makeatother
% Scale images if necessary, so that they will not overflow the page
% margins by default, and it is still possible to overwrite the defaults
% using explicit options in \includegraphics[width, height, ...]{}
\setkeys{Gin}{width=\maxwidth,height=\maxheight,keepaspectratio}
% Set default figure placement to htbp
\makeatletter
\def\fps@figure{htbp}
\makeatother
\setlength{\emergencystretch}{3em} % prevent overfull lines
\providecommand{\tightlist}{%
  \setlength{\itemsep}{0pt}\setlength{\parskip}{0pt}}
\setcounter{secnumdepth}{-\maxdimen} % remove section numbering
\usepackage{bm, amsmath, amssymb}
\ifLuaTeX
  \usepackage{selnolig}  % disable illegal ligatures
\fi
\IfFileExists{bookmark.sty}{\usepackage{bookmark}}{\usepackage{hyperref}}
\IfFileExists{xurl.sty}{\usepackage{xurl}}{} % add URL line breaks if available
\urlstyle{same}
\hypersetup{
  pdftitle={Introduction to Statistics and Computation with Data: HomeWork 8},
  pdfauthor={Daibik Barik (Bmat2316)},
  hidelinks,
  pdfcreator={LaTeX via pandoc}}

\title{\textbf{Introduction to Statistics and Computation with Data:
HomeWork 8}}
\author{Daibik Barik (Bmat2316)}
\date{2024-03-20}

\begin{document}
\maketitle

\section*{1. }

Constructing a confidence interval for the given data,

Let \(X_1, X_2, \dots, X_{n=8} = 300,500,300,700,300,500,300,700\)
respectively.
\[\overline X = \frac{300+500+300+700+300+500+300+700}{8} = \frac{3600}{8} = 450\]
It is given that each \(X_i \sim \mathcal{N}(\mu, \sigma^2)\), where
\(\sigma^2 = 100\)
\[\implies K =\frac{\sqrt{n}(\bar{X} - \mu)}{\sigma} \sim \mathcal{N}(0,1)\]
To construct a \(95\%\) confidence interval, we need to find `\(z\)'
such that
\[P\left[\mu \in \left(\bar X - \frac{z\sigma}{\sqrt{n}},\bar X + \frac{z\sigma}{\sqrt{n}}\right)\right] = 0.95\]

\[P\left[\mu \in \left(450 - \frac{z\sqrt{100}}{\sqrt{8}},450 + \frac{z\sqrt{100}}{\sqrt{8}} \right) \right] = 0.95\]

\[P\left[\frac{\sqrt{2}(\mu - 450)}{5} \in \left( - z,+ z \right) \right] = 0.95\]
\[\implies P[-z< K < z] = 0.95, \text{ where } K \sim \mathcal{N}(0,1) \]
\[\implies \Phi(z) - \Phi(-z) = 0.95\]
\[\implies \Phi(z) - (1 - \Phi(z)) = 0.95\] \[\implies \Phi(z) = 0.975\]

\begin{Shaded}
\begin{Highlighting}[]
\FunctionTok{qnorm}\NormalTok{(}\FloatTok{0.975}\NormalTok{)}
\end{Highlighting}
\end{Shaded}

\begin{verbatim}
## [1] 1.959964
\end{verbatim}

\[\implies z = 1.959964 \approx 1.96\] Hence the interval we get is
\[\left(450 - \frac{(1.96)\sqrt{100}}{\sqrt{8}},450 + \frac{(1.96)\sqrt{100}}{\sqrt{8}} \right)\]
\[ = \left( 450 - 6.93,450 + 6.93 \right)\]
\[ = \left( 443.07,456.93\right)\]

\section*{2. }

\subsection*{(a). }

\(X_1, ..., X_n\) are iid \(Uniform(0,\theta)\)
\[L(\theta| X_1, ..., X_n) = \prod_{i=1}^n\frac{1}{\theta} = \theta^{-n}\]
\[\log(L(\theta| X_1, ..., X_n)) = -n\cdot \log(\theta)\]

Differentiating on both sides,
\[\frac{d\log(L(\theta| X_1, ..., X_n))}{d\theta} = \frac{1}{L(\theta| X_1, ..., X_n)}\frac{-n}{\theta}\]
\[\implies \frac{-n}{\theta} \rightarrow 0\]

This happens when \(\theta\) takes the maximum possible value,which is
\(Max\{X_1, ..., X_n\}\)

\subsection*{(b). }

\(\hat{\theta} = Max\{X_1, ..., X_n\}\) and we have to find the
distribution of \(T = \frac{ \hat{ \theta }}{ \theta }\)

\[P\left[ \frac{\hat{\theta}}{\theta} \leq t\right] = P[\hat{\theta} \leq t\theta]\]
\[ = P[Max\{X_1, ..., X_n\} \leq t\theta]\]
\[ = P[X_1 < t\theta] \cap ... \cap P[X_n < t\theta]\]
\[ = \left(\frac{t\theta}{\theta}\right)^n\] \[ = t^n\]

\subsection*{(c). }

\section*{3. }

\[\frac{\sqrt{n}(\overline{X} - \mu)}{S} \sim t_{n-1}\]
\[K = \frac{\sqrt{19}(12 - \mu)}{6.3} \sim t_{18}\]

To construct a 95\% confidence interval for the true mean,
\[P\left[\mu \in \left(12 -\frac{z\cdot6.3}{\sqrt{19}},12 +\frac{z\cdot6.3}{\sqrt{19}} \right)\right] = 0.95\]
\[P[K \in (-z,z) = 0.95]\]

\begin{Shaded}
\begin{Highlighting}[]
\FunctionTok{qt}\NormalTok{(}\FloatTok{0.975}\NormalTok{, }\DecValTok{18}\NormalTok{)}
\end{Highlighting}
\end{Shaded}

\begin{verbatim}
## [1] 2.100922
\end{verbatim}

\(\implies z = 1.734064 \approx 1.7\)

So the 95\% confidence interval is,
\[\left(12 -\frac{2.1\times 6.3}{\sqrt{19}},12 +\frac{2.1\times6.3}{\sqrt{19}}\right)\]
\[= \left(8.96483,15.03517\right)\]

\section*{4. }

Given,
\(X_i \sim Bernoulli(p), \text{ } \forall \text{ } x \in \{1,2, .. n\}\)
We estimate p as \(\hat{p} = \frac{1}{n}\sum_1^n X_i\) For the bernoulli
distribution, the standard deviation is estimated by,
\[\sqrt{\hat{p}(1-\hat{p})}\] Now,the Central Limit Theorem says,
\[\frac{\sqrt{n}(X - \hat{p})}{\sqrt{\hat{p}(1-\hat{p})}} \sim \mathcal{N}(0,1)\]
So, to find the 95\% confidence interval we have to find z such that,
\[P\left[K = \frac{\sqrt{n}(X - \hat{p})}{\sqrt{\hat{p}(1-\hat{p})}} \in (-z,z) \right] = 0.95\]
\[\implies P[-z < K < z] = 0.95, \text{ where } K \sim N(0,1) \]
\[\implies \Phi(z) - \Phi(-z) = 0.95\] \[\implies \Phi(z) = 0.975\]

\begin{Shaded}
\begin{Highlighting}[]
\FunctionTok{qnorm}\NormalTok{(}\FloatTok{0.975}\NormalTok{)}
\end{Highlighting}
\end{Shaded}

\begin{verbatim}
## [1] 1.959964
\end{verbatim}

\[\implies z = 1.959964 \approx 1.96\] So, we say that the 95\%
confidence interval for p is
\[\left(\hat{p} - 1.96\sqrt{ \frac{\hat{p}(1-\hat{p})}{n}}, \hat{p} + 1.96\sqrt{ \frac{\hat{p}(1-\hat{p})}{n}}\right)\]

\section*{5. }
\subsection*{(a). }

Given,
\(X_i \sim \mathcal{N}(0,\sigma^2), \text{ } \forall \text{ } i \in \{1,2, .. n\}\)
\[\implies \frac{X_i}{\sigma} \sim \mathcal{N}(0,1)\]
\[\implies \left(\frac{X_i}{\sigma}\right)^2 \sim \Gamma\left(\frac{1}{2}, \frac{1}{2}\right)\]
Now, by the additivity of the \(\Gamma\) distribution,
\[\implies \left(\sum_{i=1}^n\left(\frac{X_i}{\sigma}\right)^2\right) \sim \Gamma\left(\frac{n}{2}, \frac{1}{2}\right) \sim \chi_n^2\]
Hence, \(\left(\sum_{i=1}^n\left(\frac{X_i}{\sigma}\right)\right)^2\)
follows the chi-squared distribution with n degrees of freedom.

\newpage
\subsection*{(b). }

\begin{Shaded}
\begin{Highlighting}[]
\NormalTok{alpha }\OtherTok{=} \FloatTok{0.95}
\NormalTok{n }\OtherTok{=} \DecValTok{100}
\NormalTok{lower\_value }\OtherTok{=} \FunctionTok{qchisq}\NormalTok{((}\DecValTok{1}\SpecialCharTok{{-}}\NormalTok{alpha)}\SpecialCharTok{/}\DecValTok{2}\NormalTok{,n)}
\NormalTok{upper\_value }\OtherTok{=} \FunctionTok{qchisq}\NormalTok{((}\DecValTok{1}\SpecialCharTok{+}\NormalTok{alpha)}\SpecialCharTok{/}\DecValTok{2}\NormalTok{,n)}
\FunctionTok{cat}\NormalTok{(}\StringTok{"The 95\% confidence interval for the chi{-}squared distribution }
\StringTok{            with 100 degrees of freedom is ("}\NormalTok{, lower\_value, }\StringTok{","}\NormalTok{, upper\_value , }\StringTok{")"}\NormalTok{)}
\end{Highlighting}
\end{Shaded}

\begin{verbatim}
## The 95% confidence interval for the chi-squared distribution 
##             with 100 degrees of freedom is ( 74.22193 , 129.5612 )
\end{verbatim}

\subsection{(c). }

\begin{Shaded}
\begin{Highlighting}[]
\NormalTok{alpha }\OtherTok{=} \FloatTok{0.95}
\NormalTok{n }\OtherTok{=} \DecValTok{100}
\NormalTok{lower\_value }\OtherTok{=} \FunctionTok{qchisq}\NormalTok{(}\DecValTok{1}\SpecialCharTok{{-}}\NormalTok{alpha,n)}
\NormalTok{upper\_value }\OtherTok{=} \FunctionTok{qchisq}\NormalTok{(alpha,n)}
\FunctionTok{cat}\NormalTok{(}\StringTok{"The 95\% confidence interval for the chi{-}squared distribution }
\StringTok{            with 100 degrees of freedom is either"}\NormalTok{, }\StringTok{"("}\NormalTok{, }\DecValTok{0}\NormalTok{, }\StringTok{","}\NormalTok{, upper\_value , }\StringTok{")"}\NormalTok{, }
\NormalTok{            (}\StringTok{"or"}\NormalTok{), }\StringTok{"("}\NormalTok{,lower\_value, }\StringTok{","}\NormalTok{, }\ConstantTok{Inf}\NormalTok{ , }\StringTok{")"}\NormalTok{ )}
\end{Highlighting}
\end{Shaded}

\begin{verbatim}
## The 95% confidence interval for the chi-squared distribution 
##             with 100 degrees of freedom is either ( 0 , 124.3421 ) or ( 77.92947 , Inf )
\end{verbatim}

\end{document}
